\section{Related Work}
\subsection{Legged gaze and camera stabilization}
Shaw et al. developed workspace CPG with SE(3) body-pose control on an 18-DOF hexapod, enabling online visually guided locomotion and directed gaze without an independent camera actuator \cite{shaw2019}. It is therefore incorrect to distinguish our study by claiming that earlier controllers used only inertial sensing or did not alter CPG motion using vision. The proposed distinction is the downstream residual pixel-motion criterion and its explicit dependence on the compensation limits of each gait candidate. Earlier camera-oriented gait learning also exists: Saggar et al. jointly optimized speed and visual stability on an Aibo quadruped and demonstrated a perception benefit \cite{saggar2007}. That study is a substantive precedent for image-quality-aware gait search.

Motion-blur-constrained visual control has also been formulated as an optimization problem; Robic et al. incorporated a blur-related constraint into QP visual servoing \cite{robic2024}. Thus the use of a camera-space cost or of a constrained optimization solver is an established technique. Our objective is not a new visual servo controller but a locomotion-specific choice among disturbances that can or cannot be compensated by the body.

\subsection{Hexapod gait and body coordination}
Burkus et al. explicitly minimized unwanted body motion, including the sensor consequences of vibration, while co-optimizing structural and walking parameters \cite{burkus2022}. Smart Gait addressed hexapod swing trajectories and gait tuning to reduce mechanical costs and improve locomotion behavior \cite{yin2024}. Gao et al. combined gait planning, smoothed foot trajectories and contact-related compliant control to improve motion stability \cite{gao2022}. These studies preclude claims that we are first to reduce vibration at its source by changing a walking gait.

Yu et al. coordinated body pose and center-of-gravity adjustment across complex-terrain support phases \cite{yu2025}; Liu et al. used an MPC layer to map motion decisions to CPG parameters in underwater hexapod tracking \cite{liu2026}. Both results establish mature frameworks for coupling upper-level optimization with lower-level motion generation. Our controller adopts such a framework rather than claiming it as a contribution.

\subsection{Specific methodological distinction}
Prior work overlaps with one or more components of our proposed system: gait-induced motion mitigation, camera-oriented gait optimization, body-pose compensation, and supervisory CPG adaptation. The testable distinction in this study is the \emph{candidate-conditioned achievable image motion after a support- and actuator-constrained compensator}, rather than an additive visual penalty on raw gait motion. This distinction should be demonstrated against both (a) raw-motion gait selection followed by the identical compensator and (b) a simplified candidate selector that assumes constant compensation authority. The absence of a directly identical method in a finite literature search is not by itself evidence of novelty.
